%-*-latex-*-
%-*-latex-*-
\newcommand\COURSE{ciss240}
\newcommand\ASSESSMENT{q1402}
\newcommand\ASSESSMENTTYPE{Quiz}
\newcommand\POINTS{\textwhite{xxx/xxx}}

\input{myquizpreamble}
\input{yliow}
\input{\COURSE}
%-*-latex-*-

%-*-latex-*-

%-*-latex-*-
\renewcommand\TITLE{\ASSESSMENTTYPE \ \ASSESSMENT}

\renewcommand\EMAIL{}
\renewcommand\AUTHOR{eagraham1@cougars.ccis.edu}

\textwidth=6in
\begin{document}
\topmatterthree

This is a closed-book, no compiler, 5 minute quiz.

%------------------------------------------------------------------------------
\nextq
Write a program that gets two integer $n$ and $i$ from the user
and prints $i$--th digit of $n$:
\begin{itemize}[nosep]
\li if $n = 13524$ and $i = 0$, then $4$ is printed,
\li if $n = 13524$ and $i = 1$, then $2$ is printed,
\li if $n = 13524$ and $i = 2$, then $5$ is printed,
\li if $n = 13524$ and $i = 3$, then $3$ is printed,
\li if $n = 13524$ and $i = 4$, then $1$ is printed,
\li if $n = 13524$ and $i = 5$, then $0$ is printed,
\li if $n = 13524$ and $i = 6$, then $0$ is printed,
\end{itemize}
Here's a test run:
\vspace{-5pt}
\begin{Verbatim}[fontsize=\footnotesize,commandchars=\\\{\},frame=single]
\userinput{13524 2}
5
\end{Verbatim}
\ANSWER
\begin{answercode}
#include <iostream>

int main()
{
    int n, i;
    std::cin >> n >> i;

    int digit = 0;

    if (n < 0)
    {
        n = -n;
    }

    for (int j = 0; j <= i; ++j)
    {
        digit = n % 10; 
        n /= 10;
    }

    std::cout << digit << '\n';
    return 0;
}

\end{answercode}

%------------------------------------------------------------------------------
\nextq
Write a program that gets an integer $n$ from the user and
continually
replace $n$ by the sum of the digits of $n$
until a one digit $n$ is reached.
The values taken by $n$ is printed.
For instance if the user enters $n = 123456$, then
$n = 1+2+3+4+5+6 = 21$,
$n = 2+1 = 3$.
Here's a test run:
\vspace{-5pt}
\begin{Verbatim}[fontsize=\footnotesize,commandchars=\\\{\},frame=single]
\userinput{123456}
21
3
\end{Verbatim}
\ANSWER
\begin{answercode}
#include <iostream>
int main()
{
    int n;
    std::cin >> n;

    int num_digits; // number of digits in n
    
    // compute num_digits
    num_digits = 0;
    int temp = n;
    
    while (temp > 0)
    {
        temp /= 10;
        num_digits++;
    }

    while (num_digits > 1)
    {
        int sum_digits = 0;

        // compute sum_digits
        temp = n;
        while (temp > 0)
        {
            sum_digits += temp % 10; 
            temp /= 10;              
        }

        // print sum_digits
        std::cout << sum_digits << '\n';

        // get ready for the next iteration
        n = sum_digits;

        // compute num_digits
        num_digits = 0;
        temp = n;
        while (temp > 0)
        {
            temp /= 10;
            num_digits++;
        }
    }

    return 0;
}
\end{answercode}

%------------------------------------------------------------------------------
%-*-latex-*-
\newpage
\input{instructions}
\end{document}

